% ============================================================
% 6. DISCUSSION
% ============================================================
\section{Discussion}

\subsection{Limitations}

\textbf{Reliance on LLM-as-a-Judge.} PCS and EmotVar rely on LLM-based emotion extraction (via the actor LLM itself; \S4.4). Recent work demonstrates that strong LLM judges achieve over 80\% agreement with human annotators on complex conversational metrics, matching inter-annotator agreement levels \cite{zheng2023judging}. Nevertheless, LLM-based evaluation may exhibit systematic biases (e.g., verbosity preference, self-enhancement) that could affect emotion extraction. We did not assess inter-rater reliability of the critic across repeated evaluations of identical inputs. Future work incorporating targeted human evaluation could further calibrate the correlation between our automated EmotVar and human-perceived emotional authenticity. The keyword-based Collapse Index, while deterministic and reproducible, may miss subtle template patterns or produce false positives on genuinely warm responses.

\textbf{Limited model coverage.} Three API-served models do not constitute a comprehensive alignment survey. Notably absent are open-weight models where alignment strength can be controlled.

\textbf{MiniMax data quality.} Systematic failure on the \texttt{emotional\_comfort} scenario inflates variance. Robustness checks show consistent trends but are underpowered.

\textbf{Ablation scope.} Our ablation demonstrates component necessity but not physics-specificity. A controlled comparison against equivalently complex non-physics alternatives (e.g., fixed-schedule noise, recency-weighted retrieval) would be needed to isolate the contribution of the physics formulation itself.

\textbf{Small $N$ for secondary findings.} With $N = 15$ per cell, the primary finding is well-powered, but secondary findings (ablation, MiniMax) are underpowered.

\textbf{Single persona archetype.} All experiments use a single tsundere persona. TPE inherently assumes emotional variance is desirable; for naturally stoic or reserved personas (e.g., a monk, a military commander), the temperature scaling factor may require persona-conditioned dampening to avoid artificially inflating expressiveness.

\subsection{Future Work}

Three directions emerge: (1) \emph{alignment-aware persona control} combining TPE with model-specific fine-tuning or contrastive decoding to bypass the alignment ceiling; (2) \emph{physics-specificity ablation} comparing against fixed-variance noise and recency-weighted retrieval; and (3) \emph{multi-persona evaluation} extending to diverse character archetypes, including low-variance personas where TPE's temperature coupling would need adaptation.

% ============================================================
% 7. CONCLUSION
% ============================================================
\section{Conclusion}

We presented the Thermodynamic Persona Engine (TPE), a physics-inspired framework that couples frustration-driven temperature to behavioral signal noise and uses mass-weighted gravitational memory with temporal decay to anchor response style in LLM role-playing agents. On qwen3-max, TPE achieves a 32\% improvement in emotional variance without degrading persona consistency (Cliff's $d = +0.75$, large effect). Beyond TPE itself, our cross-model evaluation surfaces two findings of broader significance: Reflexion-style self-correction---widely adopted in agent architectures---systematically degrades persona consistency, and strong RLHF alignment creates a hard ceiling on prompt-level persona control, with anti-alignment instructions paradoxically triggering the very safety-trained patterns they aim to suppress. These results suggest that emotionally rich character simulation requires not only better prompting frameworks but also alignment-aware approaches that work \emph{with} rather than against the model's post-training distribution.

% ============================================================
% 8. ETHICAL CONSIDERATIONS
% ============================================================
\section{Ethical Considerations}

Our research explores methods to modulate RLHF-aligned behavior for realistic persona simulation, which raises dual-use concerns. We explicitly acknowledge the following:

\textbf{Intended scope.} TPE is designed for controlled creative environments---interactive fiction, game NPCs, social skill training simulations, and research on LLM expressiveness. It should \emph{not} be deployed in open-domain public chatbots, therapeutic applications, or customer service systems where empathetic alignment is critical for user safety.

\textbf{Anti-alignment instructions.} Our system prompt instructs the actor LLM to prioritize character-consistent expression over default assistant behavior. While this is necessary for persona fidelity, the same technique could be misused to suppress safety-relevant responses. Our gpt-5-mini results (\S\ref{sec:alignment}) demonstrate that strongly aligned models resist such override attempts, suggesting that current safety training provides meaningful protection against prompt-level anti-alignment.

\textbf{Simulated emotions.} The ``frustration,'' ``temperature,'' and ``drives'' in TPE are computational abstractions operating on scalar vectors. The agent does not experience sentience, suffering, or genuine emotional states. We caution against anthropomorphizing these mechanisms.
